% !TeX root = ../教材题库_全部条目_转写工作稿.tex
% Source: XJ-B1-C06-HW-P232-Y0932
% 第6章 / 6.3 统计图表 / 教材 p.232、233 / PDF 第239、240页
% OCR未校对：公式与文字均可能有误；未完成数学TeX转写；禁止直接发给学生
% 原题及相邻原解截图（核对用）：../题目原图/XJ-B1-C06-HW-P232-Y0932-01.webp, ../题目原图/XJ-B1-C06-HW-P232-Y0932-02.webp
% 本文件只在转写完成后作为题干使用；OCR 条目尚不是可靠题干。
\begin{bankitem}{习题6.3 第5题}{第6章 / 6.3 统计图表 / 教材 p.232、233 / PDF 第239、240页}{XJ-B1-C06-HW-P232-Y0932}{OCR未校对：公式与文字均可能有误；未完成数学TeX转写；禁止直接发给学生}
\begin{ocrdraft}
[教材p.232 / PDF第239页]
\par
5．为了解游客人数的变化规律，提高旅游服务质量，一个研究组收集并整理
\par
了所在城市连续三年的月接待游客量的数据，绘制了下面的折线图．
\par
月接待游客量/万人次
\par
45
\par
40
\par
35
\par
25
\par
12345
\par
6
\par
7
\par
8
\par
9]0ll12]
\par
2
\par
3
\par
45678
\par
910111212345678910H]2
\par
第一年
\par
第二年
\par
第三年
\par
（第5题）
\par
根据该折线图，下列结论正确的有
\par
(A)月接待游客量逐月增加
\par
（B)年接待游客量逐年增加
\par
（0各年的月接待游客量高峰期大致在7一8月
\par
(D)各年1月至6月的月接待游客量相对于7月至12月，波动性更小，变化
\par
[教材p.233 / PDF第240页]
\par
比较平稳
\end{ocrdraft}
\end{bankitem}
